\section{One level of Hastings' hierarchical projection}
\label{sec:hastings05_hierarchical_projection}

% Hastings, arXiv:cond-mat/0508554v3, proj3new.tex, lines 596--756.

Section~II of~\cite{Hastings2006Solving} approximates a gapped ground state
$\Psi_a$ by hierarchical spectral projections of pair operators $M=M_{k,k+1}$
satisfying $\lVert(M-c)\Psi_a\rVert\le\epsilon$ for
$c=\langle\Psi_a,M\Psi_a\rangle$. Level~$1$ projects onto eigenvalues of $M$ in
$[c-x,c+x]$; each subsequent level compresses $M$ to the range of the previous
level and normalizes the projected state over $h=\log_2N$ levels. Here level
deficiencies are measured in norm rather than expectation so that the
next-level normalization bound holds
\cite{gap:hastings05_hierarchical_projection}. The full $h$-level hierarchy is
in Section~\ref{sec:hastings05_hierarchy_levels}.

\begin{definition}[Window projection]
    \label{def:hastings05_window_projection}
    \lean{TNLean.MPS.AreaLaw.LocalSolution.windowProjection}
    \leanok
    \uses{def:hastings_spectral_cutoff}
    For $M=M^\dagger$ and $c,x\in\mathbb R$, the \emph{window projection}
    $P_x(M,c)=1_{[-x,x]}(M-c)$ is the spectral projection of $M$ onto
    $[c-x,c+x]$ ($P^1_k$ of~\cite[lines 640--642]{Hastings2006Solving} at
    $c=\langle\Psi_a,M\Psi_a\rangle$).
\end{definition}

\begin{theorem}[The shifted operator]
    \label{thm:hastings05_shifted_hermitian}
    \lean{TNLean.MPS.AreaLaw.LocalSolution.isHermitian_sub_ofReal_smul_one,
        TNLean.MPS.AreaLaw.LocalSolution.ofReal_smul_one_eq_algebraMap,
        TNLean.MPS.AreaLaw.LocalSolution.isStarProjection_windowProjection}
    \leanok
    \uses{def:hastings05_window_projection,thm:hastings_spectral_cutoff_projection}
    For $M=M^\dagger$ and $c,x\in\mathbb R$, $M-c$ is Hermitian and $P_x(M,c)$
    is an orthogonal projection.
\end{theorem}
\begin{proof}\leanok
    Apply Theorem~\ref{thm:hastings_spectral_cutoff_projection} to
    $(M-c)^\dagger=M-c$.
\end{proof}

\begin{theorem}[Deficiency and expectation of a projection]
    \label{thm:hastings05_projection_deficiency}
    \lean{TNLean.MPS.AreaLaw.LocalSolution.re_inner_toEuclideanLin_eq_norm_sq,
        TNLean.MPS.AreaLaw.LocalSolution.norm_sq_add_norm_sub_sq,
        TNLean.MPS.AreaLaw.LocalSolution.re_inner_eq_one_sub_norm_sub_sq,
        TNLean.MPS.AreaLaw.LocalSolution.one_sub_le_re_inner_of_norm_sub_le,
        TNLean.MPS.AreaLaw.LocalSolution.one_sub_sq_le_re_inner_of_norm_sub_le,
        TNLean.MPS.AreaLaw.LocalSolution.norm_sub_mul_sq_le_of_isStarProjection}
    \leanok
    For an orthogonal projection $P$ and $\|\Psi\|=1$,
    $\langle\Psi,P\Psi\rangle=\lVert P\Psi\rVert^2=1-\lVert\Psi-P\Psi\rVert^2$;
    if $\lVert\Psi-P\Psi\rVert\le\delta$, then
    $\langle\Psi,P\Psi\rangle\ge1-\delta^2\ge1-\delta$. For any matrix $P'$ and
    vector $\Psi$,
    \begin{align}
        \lVert\Psi-PP'\Psi\rVert^2
         & \le\lVert\Psi-P\Psi\rVert^2+\lVert\Psi-P'\Psi\rVert^2.
        \label{eq:hastings05_two_projections_sq}
    \end{align}
\end{theorem}
\begin{proof}\leanok
    \uses{thm:hastings_final_state_bessel}
    By Theorem~\ref{thm:hastings_final_state_bessel} and orthogonality of
    $P\Psi$ and $(1-P)\Psi$, $\langle\Psi,P\Psi\rangle=\|P\Psi\|^2=1-d^2$ with
    $d=\|\Psi-P\Psi\|\le1$, so $d^2\le d\le\delta$; similarly
    $\Psi-PP'\Psi=(1-P)\Psi+P(1-P')\Psi$ gives
    (\ref{eq:hastings05_two_projections_sq}).
\end{proof}

\begin{theorem}[One level]
    \label{thm:hastings05_one_level}
    \lean{TNLean.MPS.AreaLaw.LocalSolution.norm_sub_windowProjection_le,
        TNLean.MPS.AreaLaw.LocalSolution.one_sub_div_le_re_inner_windowProjection,
        TNLean.MPS.AreaLaw.LocalSolution.one_sub_div_sq_le_re_inner_windowProjection}
    \leanok
    \uses{def:hastings05_window_projection,thm:hastings05_shifted_hermitian,
        thm:hastings05_projection_deficiency,thm:hastings_spectral_cutoff_obd}
    For $M=M^\dagger$, $c\in\mathbb R$, $x>0$, $P=P_x(M,c)$, and any $\Psi$,
    \begin{align}
        \lVert\Psi-P\Psi\rVert
         & \le\frac{\lVert(M-c)\Psi\rVert}{x},
        \label{eq:hastings05_one_level}
    \end{align}
    and if $\|\Psi\|=1$, then
    $\langle\Psi,P\Psi\rangle\ge1-(\|(M-c)\Psi\|/x)^2\ge1-\|(M-c)\Psi\|/x$
    \cite[lines 645--657]{Hastings2006Solving}.
\end{theorem}
\begin{proof}\leanok
    Combine Theorem~\ref{thm:hastings_spectral_cutoff_obd} for $M-c$ with
    Theorem~\ref{thm:hastings05_projection_deficiency}.
\end{proof}

\begin{theorem}[An eigenvalue in the window]
    \label{thm:hastings05_eigenvalue_in_window}
    \lean{TNLean.MPS.AreaLaw.LocalSolution.exists_mem_spectrum_abs_sub_le,
        TNLean.MPS.AreaLaw.LocalSolution.windowProjection_eq_zero}
    \leanok
    \uses{thm:hastings05_one_level}
    If $M=M^\dagger$, $c\in\mathbb R$, $x>0$, and $0<\|\Psi\|$ with
    $\|(M-c)\Psi\|<x\|\Psi\|$, then $\operatorname{spec}(M)\cap[c-x,c+x]\ne\emptyset$
    \cite[lines 658--667]{Hastings2006Solving}.
\end{theorem}
\begin{proof}\leanok
    If $P_x(M,c)=0$, (\ref{eq:hastings05_one_level}) gives the contradiction
    $\|\Psi\|\le\|(M-c)\Psi\|/x<\|\Psi\|$.
\end{proof}

\begin{theorem}[Normalization]
    \label{thm:hastings05_normalization}
    \lean{TNLean.MPS.AreaLaw.LocalSolution.norm_inv_norm_smul_sub_le,
        TNLean.MPS.AreaLaw.LocalSolution.norm_inv_norm_smul_toEuclideanLin_sub_sq_le,
        TNLean.MPS.AreaLaw.LocalSolution.norm_inv_norm_smul_sub_le_two_mul}
    \leanok
    \uses{thm:hastings05_projection_deficiency}
    Let $\|\Psi\|=1$ (with $0/0=0$). For any $v$,
    $\|v/\|v\|-\Psi\|\le2\|v-\Psi\|$; for an orthogonal projection $Q$,
    $\|Q\Psi/\|Q\Psi\|-\Psi\|^2\le2\|\Psi-Q\Psi\|^2$; and for commuting
    orthogonal projections $P,P'$ with $\|\Psi-P\Psi\|,\|\Psi-P'\Psi\|\le b$,
    \begin{align}
        \Bigl\lVert\frac{PP'\Psi}{\lVert PP'\Psi\rVert}-\Psi\Bigr\rVert
         & \le2b
        \label{eq:hastings05_normalization}
    \end{align}
    \cite[lines 686--689]{Hastings2006Solving}.
\end{theorem}
\begin{proof}\leanok
    Use $|1-\|v\||\le\|v-\Psi\|$, $2(1-t)\le2(1-t^2)$ for
    $t=\|Q\Psi\|\in(0,1]$ ($1\le2$ if $Q\Psi=0$), and
    (\ref{eq:hastings05_two_projections_sq}) for $Q=PP'$.
\end{proof}

\begin{theorem}[Nested projections]
    \label{thm:hastings05_nested_projections}
    \lean{TNLean.MPS.AreaLaw.LocalSolution.norm_sub_sq_le_of_mul_eq}
    \leanok
    \uses{thm:hastings05_projection_deficiency}
    For orthogonal projections $P,Q$ with $PQ=P$, $\|\Psi\|=1$, and
    $\Psi'=Q\Psi/\|Q\Psi\|$,
    \begin{align}
        \lVert\Psi-P\Psi\rVert^2
         & \le\lVert\Psi-Q\Psi\rVert^2+\lVert\Psi'-P\Psi'\rVert^2
        \label{eq:hastings05_nested_projections}
    \end{align}
    \cite[lines 671--676]{Hastings2006Solving}.
\end{theorem}
\begin{proof}\leanok
    With $t=\|Q\Psi\|\le1$ and $s=\|P\Psi'\|\le1$, $\|P\Psi\|=ts$ and
    $1-t^2s^2=(1-t^2)+t^2(1-s^2)\le(1-t^2)+(1-s^2)$.
\end{proof}

\begin{theorem}[The residual of a nearby vector]
    \label{thm:hastings05_residual}
    \lean{TNLean.MPS.AreaLaw.LocalSolution.norm_toEuclideanLin_le_add,
        TNLean.MPS.AreaLaw.LocalSolution.norm_sub_ofReal_smul_one_le,
        TNLean.MPS.AreaLaw.LocalSolution.norm_sub_smul_one_le_two_mul,
        TNLean.MPS.AreaLaw.LocalSolution.abs_re_inner_toEuclideanLin_le}
    \leanok
    For any $N,\phi,\psi$,
    $\|N\phi\|\le\|N\psi\|+\|N\|\,\|\phi-\psi\|$; for $M$ and $c\in\mathbb R$,
    $\|M-c\|\le\|M\|+|c|$ (and $\|M-c\|\le2\|M\|$ when $|c|\le\|M\|$, which
    holds for $c=\operatorname{Re}\langle\psi,M\psi\rangle$ when $\|\psi\|=1$)
    \cite[lines 690--695]{Hastings2006Solving}.
\end{theorem}
\begin{proof}\leanok
    Apply the triangle inequality to $N\phi=N\psi+N(\phi-\psi)$ and $M-c$, and
    Cauchy--Schwarz to $\langle\psi,M\psi\rangle$.
\end{proof}

\begin{theorem}[One level combined]
    \label{thm:hastings05_one_level_combined}
    \lean{TNLean.MPS.AreaLaw.LocalSolution.norm_sub_windowProjection_le_of_norm_sub_le}
    \leanok
    \uses{thm:hastings05_one_level,thm:hastings05_residual}
    If $M=M^\dagger$, $c\in\mathbb R$, $x>0$, $\|(M-c)\psi\|\le\epsilon$, and
    $\|\phi-\psi\|\le\eta$, then
    \begin{align}
        \lVert\phi-P_x(M,c)\phi\rVert
         & \le\frac{\epsilon+\lVert M-c\rVert\,\eta}{x}
        \label{eq:hastings05_one_level_combined}
    \end{align}
    \cite[lines 690--699]{Hastings2006Solving}.
\end{theorem}
\begin{proof}\leanok
    Combine (\ref{eq:hastings05_one_level}) for $\phi$ with
    Theorem~\ref{thm:hastings05_residual}.
\end{proof}

\begin{theorem}[The residual under a compression]
    \label{thm:hastings05_compression}
    \lean{TNLean.MPS.AreaLaw.LocalSolution.norm_toEuclideanLin_compress_sub_le}
    \leanok
    For an orthogonal projection $Q$, matrix $M$, $c\in\mathbb R$, and
    $Q\phi=\phi$, $\|(QMQ-c)\phi\|\le\|(M-c)\phi\|$
    \cite[lines 673--689]{Hastings2006Solving}.
\end{theorem}
\begin{proof}\leanok
    Since $Q\phi=\phi$ and $\|Q\|\le1$,
    $\|(QMQ-c)\phi\|=\|Q(M-c)\phi\|\le\|(M-c)\phi\|$.
\end{proof}

\begin{definition}[Window projection in the preceding subspace]
    \label{def:hastings05_compressed_window}
    \lean{TNLean.MPS.AreaLaw.LocalSolution.compressedWindowProjection}
    \leanok
    \uses{def:hastings05_window_projection}
    For an orthogonal projection $Q$, $M=M^\dagger$, and $c,x\in\mathbb R$, set
    $P=QP_x(QMQ,c)$ \cite[lines 673--685]{Hastings2006Solving}.
\end{definition}

\begin{theorem}[Restriction to the preceding subspace]
    \label{thm:hastings05_compressed_window_projection}
    \lean{TNLean.MPS.AreaLaw.LocalSolution.commute_windowProjection_compress,
        TNLean.MPS.AreaLaw.LocalSolution.isStarProjection_compressedWindowProjection,
        TNLean.MPS.AreaLaw.LocalSolution.compressedWindowProjection_mul,
        TNLean.MPS.AreaLaw.LocalSolution.toEuclideanLin_compressedWindowProjection_of_mem,
        TNLean.MPS.AreaLaw.LocalSolution.toEuclideanLin_compressedWindowProjection_eq_zero}
    \leanok
    \uses{def:hastings05_compressed_window,thm:hastings05_shifted_hermitian}
    In Definition~\ref{def:hastings05_compressed_window}, $W=P_x(QMQ,c)$
    commutes with $Q$, $P=QW$ is an orthogonal projection with $PQ=P$,
    $P\phi=W\phi$ when $Q\phi=\phi$, and $P\phi=0$ when $Q\phi=0$.
\end{theorem}
\begin{proof}\leanok
    Since $(QMQ)Q=Q(QMQ)$, functional calculus gives $WQ=QW$, so $P=QW$ is an
    orthogonal projection with $PQ=P$ and $P\phi=WQ\phi$.
\end{proof}

\begin{theorem}[One level inside the preceding subspace]
    \label{thm:hastings05_compressed_window_bound}
    \lean{TNLean.MPS.AreaLaw.LocalSolution.norm_sub_compressedWindowProjection_le,
        TNLean.MPS.AreaLaw.LocalSolution.norm_sub_compressedWindowProjection_le_of_norm_sub_le}
    \leanok
    \uses{thm:hastings05_compressed_window_projection,thm:hastings05_one_level,
        thm:hastings05_compression,thm:hastings05_residual}
    For an orthogonal projection $Q$, $M=M^\dagger$, $c\in\mathbb R$, $x>0$,
    $P=QP_x(QMQ,c)$, and $Q\phi=\phi$,
    \begin{align}
        \lVert\phi-P\phi\rVert
         & \le \frac{\lVert(M-c)\phi\rVert}{x}.
        \label{eq:hastings05_compressed_window_bound_1}
    \end{align}
    If also $\|(M-c)\psi\|\le\epsilon$ and $\|\phi-\psi\|\le\eta$, then
    \begin{align}
        \lVert\phi-P\phi\rVert
         & \le\frac{\epsilon+\lVert M-c\rVert\,\eta}{x}
        \label{eq:hastings05_compressed_window_bound_2}
    \end{align}
    \cite[lines 673--699]{Hastings2006Solving}.
\end{theorem}
\begin{proof}\leanok
    Since $P\phi=P_x(QMQ,c)\phi$, combine (\ref{eq:hastings05_one_level}) for
    $QMQ$ with Theorems~\ref{thm:hastings05_compression}
    and~\ref{thm:hastings05_residual}.
\end{proof}

\begin{definition}[The scalar recurrence]
    \label{def:hastings05_level_deficiency}
    \lean{TNLean.MPS.AreaLaw.LocalSolution.levelDeficiency,
        TNLean.MPS.AreaLaw.LocalSolution.levelDeficiency_zero,
        TNLean.MPS.AreaLaw.LocalSolution.levelDeficiency_succ}
    \leanok
    For $\epsilon,a,s\in\mathbb R$, define $b_0=0$ and
    $b_{m+1}=(\epsilon+ab_m)^s$ ($a=2$, $s=1-1/h$
    in~\cite[lines 696--704]{Hastings2006Solving}).
\end{definition}

\begin{theorem}[Nonnegativity of the recurrence]
    \label{thm:hastings05_level_deficiency_nonneg}
    \lean{TNLean.MPS.AreaLaw.LocalSolution.levelDeficiency_nonneg}
    \leanok
    \uses{def:hastings05_level_deficiency}
    If $\epsilon,a\ge0$ and $s\in\mathbb R$, then $b_m\ge0$ for all $m$.
\end{theorem}
\begin{proof}\leanok
    Induct on $m$ from $b_0=0$.
\end{proof}

\begin{theorem}[The recursion]
    \label{thm:hastings05_recursion}
    \lean{TNLean.MPS.AreaLaw.LocalSolution.levelDeficiency_le,
        TNLean.MPS.AreaLaw.LocalSolution.levelDeficiency_two_le}
    \leanok
    \uses{def:hastings05_level_deficiency,thm:hastings05_level_deficiency_nonneg}
    If $0<\epsilon\le1$, $a\ge0$, and $0\le s\le1$, then for every $m$,
    \begin{align}
        b_m & \le\Bigl(\sum_{j<m}a^j\Bigr)\epsilon^{s^m},
        \label{eq:hastings05_recursion}
    \end{align}
    and $b_m\le(2^m-1)\epsilon^{s^m}$ for $a=2$.
\end{theorem}
\begin{proof}\leanok
    By induction on $m$: with $C=\sum_{j<m}a^j$, $\epsilon\le\epsilon^{s^m}$ and
    $(1+aC)^s\le1+aC$ give $b_{m+1}\le(1+aC)\epsilon^{s^{m+1}}$.
\end{proof}

\begin{theorem}[The exponent]
    \label{thm:hastings05_exponent}
    \lean{TNLean.MPS.AreaLaw.LocalSolution.one_sub_inv_natCast_mem_Icc,
        TNLean.MPS.AreaLaw.LocalSolution.quarter_le_one_sub_inv_pow}
    \leanok
    For $h\in\mathbb N$ (with $1/0=0$), $0\le1-1/h\le1$, and $(1-1/h)^h\ge1/4$
    for $h\ge2$.
\end{theorem}
\begin{proof}\leanok
    Check $h=2,3$ directly; for $h=n+1\ge4$,
    $(1-1/h)^n\ge e^{-1}\ge1/3$ and $1-1/h\ge3/4$.
\end{proof}

\begin{theorem}[The bound after $h$ iterations]
    \label{thm:hastings05_final_bound}
    \lean{TNLean.MPS.AreaLaw.LocalSolution.levelDeficiency_self_le,
        TNLean.MPS.AreaLaw.LocalSolution.levelDeficiency_exp_le,
        TNLean.MPS.AreaLaw.LocalSolution.levelDeficiency_le_rpow_quarter}
    \leanok
    \uses{thm:hastings05_recursion,thm:hastings05_exponent}
    For $0<\epsilon\le1$, $h\in\mathbb N$, $a=2$, and $s=1-1/h$,
    $b_h\le(2^h-1)\epsilon^{(1-1/h)^h}$ (and $b_h\le(2^h-1)\epsilon^{1/4}$ for
    $h\ge2$). For $\epsilon=e^{-t}$ ($t\ge0$),
    \begin{align}
        b_h & \le2^h\exp\bigl(-t(1-1/h)^h\bigr)
        \label{eq:hastings05_final_bound}
    \end{align}
    \cite[lines 705--707]{Hastings2006Solving},
    \cite{gap:hastings05_hierarchical_projection}.
\end{theorem}
\begin{proof}\leanok
    Combine Theorems~\ref{thm:hastings05_recursion}
    and~\ref{thm:hastings05_exponent}.
\end{proof}
